\documentclass[12pt]{article}

\usepackage{fullpage}
\usepackage{multicol,multirow}
\usepackage{tabularx}
\usepackage{ulem}
\usepackage[utf8]{inputenc}
\usepackage{mathtools}
\usepackage{pgfplots}
\usepackage[russian]{babel}
\usepackage{fvextra}



\begin{document}

\section*{Лабораторная работа №\,3 по курсу дискрeтного анализа: исследование качества программ}

Выполнил студент группы М8О-208Б-22 МАИ \textit{Немкова Анастасия}.

\subsection*{Условие}
Для реализации словаря из предыдущей лабораторной работы, необходимо провести исследование скорости выполнения и потребления оперативной памяти. В случае выявления ошибок или явных недочётов, требуется их исправить.

Результатом лабораторной работы является отчёт, состоящий из:

\begin{itemize}
    \item Дневника выполнения работы, в котором отражено что и когда делалось, какие средства использовались и какие результаты были достигнуты на каждом шаге выполнения лабораторной работы.
    \item Выводов о найденных недочётах.
    \item Сравнение работы исправленной программы с предыдущей версией.
    \item Общих выводов о выполнении лабораторной работы, полученном опыте.
\end{itemize}

Минимальный набор используемых средств должен содержать утилиту gprof и библиотеку dmalloc, однако их можно заменять на любые другие аналогичные или более развитые утилиты (например, Valgrind или Shark) или добавлять к ним новые (например, gcov).


\subsection*{Метод решения}

Для анализа качества программ я использовала две утилиты: valgrind и gprof.

\textbf{Valgrind} - это инструментальное программоное обеспечение, предназначенное для отладки использования памяти, обнаружения утечек памяти и профилирования. 

При запуске Valgrind используются различные флаги, такие как --leak-check=full, который проверяет утечки памяти, --track-origins, который позволяет отслеживать происхождение ошибок, и --tool=memcheck, который запускает основной инструмент Memcheck для обнаружения проблем памяти.

\textbf{Gprof} - это утилита профилирования для анализа производительности программ в UNIX-подобных операционных системах. Она позволяет определить, сколько времени программа проводит в каждой функции и как часто эти функции вызываются.


\subsection*{Описание программы}

Для запуска утилиты Valgrind скомпилируем программу и пропишем в консоли valgrind --tool=memcheck --leak-check=full ./a.out

\begin{small}
\footnotesize
\begin{verbatim}
arnemkova@LAPTOP-TA2RV74U:~/DA_labs$ valgrind --tool=memcheck --leak-check=full ./a.out
==100== Memcheck, a memory error detector
==100== Copyright (C) 2002-2017, and GNU GPL'd, by Julian Seward et al.
==100== Using Valgrind-3.15.0 and LibVEX; rerun with -h for copyright info
==100== Command: ./a.out
==100==
==100==
==100== HEAP SUMMARY:
==100==     in use at exit: 122,880 bytes in 6 blocks
==100==   total heap usage: 1,609,393 allocs, 1,609,387 frees, 411,797,786 bytes allocated
==100==
==100== LEAK SUMMARY:
==100==    definitely lost: 0 bytes in 0 blocks
==100==    indirectly lost: 0 bytes in 0 blocks
==100==      possibly lost: 0 bytes in 0 blocks
==100==    still reachable: 122,880 bytes in 6 blocks
==100==         suppressed: 0 bytes in 0 blocks
==100== Reachable blocks (those to which a pointer was found) are not shown.
==100== To see them, rerun with: --leak-check=full --show-leak-kinds=all
==100==
==100== For lists of detected and suppressed errors, rerun with: -s
==100== ERROR SUMMARY: 0 errors from 0 contexts (suppressed: 0 from 0)
\end{verbatim}
\end{small}

В данном выводе HEAP SUMMARY показывает, что всего 122,880 байт в 6 блоках были использованы программой при её завершении. Общий объем памяти, выделенной и освобожденной программой, составляет 411,797,786 байт.

LEAK SUMMARY показывает, что не было обнаружено утечек памяти в трех категориях: "definitely lost" (определенно потерянная), "indirectly lost" (косвенно потерянная) и "possibly lost" (возможно потерянная). Однако было обнаружено, что 122,880 байт в 6 блоках все еще доступны (still reachable). Использование флага --show-leak-kinds=all показало, что это происходит из-за функции std::ios{\_}base::sync{\_}with{\_}stdio(bool), которая отключает синхронизацию iostream с stdio.

ERROR SUMMARY показывает, что не было обнаружено ошибок памяти в программе.\\

Для запуска утилиты Gprof:

\begin{itemize}
    \item Компилируем программу с флагом профилирования -pg
    \item Исполняем программу для создания файла с профилировочной информацией gmon.out
    \item Запускаем gprof для анализа созданного файла
\end{itemize}

Диагностика проводилась на тесте из 100- строк с операциями вставки, удаления и поиска

\begin{Verbatim}[fontsize=\footnotesize,breaklines=true]
arnemkova@LAPTOP-TA2RV74U:~/DA_labs$ g++ -pg lab2.cpp
arnemkova@LAPTOP-TA2RV74U:~/DA_labs$ ./a.out
arnemkova@LAPTOP-TA2RV74U:~/DA_labs$ gprof a.out -p gmon.out
Flat profile:

Each sample counts as 0.01 seconds.
  %   cumulative   self              self     total
 time   seconds   seconds    calls  us/call  us/call  name
 95.51      0.82     0.82  1529384     0.54     0.54  TString::TString(TString const&)
  2.33      0.84     0.02    40001     0.50     0.50  TString::operator=(TString const&)
  1.16      0.85     0.01    30001     0.33     0.33  TString::TString(char const*)
  1.16      0.86     0.01    30000     0.33     0.33  RandString(char*)
  0.00      0.86     0.00  1569387     0.00     0.00  TString::~TString()
  0.00      0.86     0.00   520108     0.00     0.00  TAVL<TString, unsigned long long>::GetHeight(TAVL<TString, unsigned long long>::TNode*)
  0.00      0.86     0.00   513315     0.00     1.07  operator>(TString const&, TString const&)
  0.00      0.86     0.00   251377     0.00     1.07  operator<(TString const&, TString const&)
  0.00      0.86     0.00   134589     0.00     0.00  TAVL<TString, unsigned long long>::FixHeight(TAVL<TString, unsigned long long>::TNode*)
  0.00      0.86     0.00   125465     0.00     0.00  TAVL<TString, unsigned long long>::BFactor(TAVL<TString, unsigned long long>::TNode*)
  0.00      0.86     0.00   120873     0.00     0.00  TAVL<TString, unsigned long long>::Balance(TAVL<TString, unsigned long long>::TNode*)
  0.00      0.86     0.00    30000     0.00    20.93  TAVL<TString, unsigned long long>::FindNode(TAVL<TString, unsigned long long>::TNode*, TString const&)
  0.00      0.86     0.00    10002     0.00     0.00  TString::TString()
  0.00      0.86     0.00    10000     0.00    19.35  TAVL<TString, unsigned long long>::InsertNode(TAVL<TString, unsigned long long>::TNode*, TAVL<TString, unsigned long long>::TNode*)
  0.00      0.86     0.00    10000     0.00    20.93  TAVL<TString, unsigned long long>::Find(TString const&)
  0.00      0.86     0.00    10000     0.00     0.50  TAVL<TString, unsigned long long>::TNode::TNode(TString const&, unsigned long long const&)
  0.00      0.86     0.00    10000     0.00     0.00  TAVL<TString, unsigned long long>::TNode::~TNode()
  0.00      0.86     0.00    10000     0.00    40.78  TAVL<TString, unsigned long long>::Insert(TString const&, unsigned long long const&)
  0.00      0.86     0.00    10000     0.00    20.93  TAVL<TString, unsigned long long>::Remove(TString const&)
  0.00      0.86     0.00     3448     0.00     0.00  TAVL<TString, unsigned long long>::RotateRight(TAVL<TString, unsigned long long>::TNode*)
  0.00      0.86     0.00     3410     0.00     0.00  TAVL<TString, unsigned long long>::RotateLeft(TAVL<TString, unsigned long long>::TNode*)
  0.00      0.86     0.00        1     0.00     0.00  _GLOBAL__sub_I__ZN7TStringC2Ev
  0.00      0.86     0.00        1     0.00     0.00  __static_initialization_and_destruction_0(int, int)
  0.00      0.86     0.00        1     0.00     0.00  TAVL<TString, unsigned long long>::Delete(TAVL<TString, unsigned long long>::TNode*)
  0.00      0.86     0.00        1     0.00     0.00  TAVL<TString, unsigned long long>::TAVL()
  0.00      0.86     0.00        1     0.00     0.00  TAVL<TString, unsigned long long>::~TAVL()
\end{Verbatim}

Данный вывод показывает, что программа большую часть времени тратит на обработку ключей, а не на операции с деревом.
95.51{\%} времени работы программы затрачивается на конструирование объектов типа TString. Далее, 2.33{\%} времени занимает оператор присваивания TString::operator=(TString const&).

В целом, чтобы улучшить производительность, следует оптимизировать работу с ключами и операции копирования и конструирования строк типа TString.


\subsection*{Выводы}

В ходе данной лабораторной работы я изучила средства диагностики и профилирования программ. В частности, я познакомилась с утилитами valgrind и gprof, которые позволяют выявлять утечки памяти, анализировать использование ресурсов и оптимизировать производительность программы.

 Эти утилиты играют важную роль в разработке программ, помогая оптимизировать код с целью уменьшения потребляемой памяти и времени работы программы.

\end{document}


